% state/state_taxonomy.tex

\chapter{State Taxonomy}

\label{chap:state-taxonomy}

Within AIM, the term \emph{state} is used at several interpretation
layers. Without an explicit taxonomy, geometric reconstruction,
realization, runtime evaluation, vehicle interpretation, and
optimization can be conflated. This chapter defines a canonical
classification that keeps these layers distinct.

\begin{principle}[State-layer principle]
\label{princ:state-layer-principle}
Within AIM, state concepts belong to distinct interpretation layers and
must not be identified implicitly.
\end{principle}

% ============================================================
\section{General State Interpretation}
% ============================================================

\begin{definition}[State]
\label{def:taxonomy-general-state}
A state is a structured description of a system configuration relative
to a specified interpretation layer.
\end{definition}

The same physical railway system can therefore induce different states,
for example geometric, realized, runtime, operational, dynamic,
vehicle-relative, or optimization states.

% ============================================================
\section{Intrinsic Geometric States}
% ============================================================

\subsection{pose2}

\begin{definition}[pose2 state]
\label{def:taxonomy-pose2}
The pose2 state is the intrinsic planar reconstruction state generated
from:
\[
\kappa(s).
\]
\end{definition}

pose2 belongs to the intrinsic geometric layer. It represents planar
geometry and tangent evolution, and it serves as the persistent
reconstruction kernel.

% ------------------------------------------------------------

\subsection{pose3}

\begin{definition}[pose3 state]
\label{def:taxonomy-pose3}
The pose3 state is the runtime-evaluated spatial railway state generated
from:
\begin{itemize}
\item pose2 reconstruction,
\item profile evaluation,
\item cant interpretation,
\item frame evaluation,
\item and runtime coordinate interpretation.
\end{itemize}
\end{definition}

pose3 describes the railway runtime state rather than the vehicle
itself.

% ============================================================
\section{Runtime and Realized States}
% ============================================================

\begin{definition}[Runtime state]
\label{def:taxonomy-runtime-state}
A runtime state is a dynamically evaluated state generated during
runtime interpretation.
\end{definition}

Runtime states depend on runtime operators, realization, coordinate
interpretation, frame construction, and operational context.

\begin{definition}[Realized state]
\label{def:taxonomy-realized-state}
A realized state is a runtime state generated from physically realized
rather than purely analytical geometry.
\end{definition}

Typical realized quantities include
\[
\mathrm{pose3}_{\mathrm{real}},
\qquad
\kappa_{\mathrm{real}},
\qquad
\phi_{\mathrm{real}}.
\]

% ============================================================
\section{Operational, Vehicle, and Dynamic States}
% ============================================================

\begin{definition}[Operational state]
\label{def:taxonomy-operational-state}
An operational state is a runtime state relevant for railway operation,
comfort, admissibility, or dynamic interpretation.
\end{definition}

Operational states are derived from runtime railway states and include,
for example, effective acceleration, jerk, roll evolution, cant
deficiency, and balance quantities.

\begin{definition}[Vehicle state]
\label{def:taxonomy-vehicle-state}
A vehicle state is a runtime-relative operational state describing the
vehicle relative to the railway runtime state.
\end{definition}

Thus,
\[
\mathrm{vehicleState}
=
\mathcal V(
\mathrm{pose3},
v,
\dots
).
\]

\begin{definition}[Center-of-gravity state]
\label{def:taxonomy-cg-state}
A center-of-gravity state is an operational abstraction describing the
runtime trajectory of the vehicle center of gravity relative to the
railway runtime state.
\end{definition}

\begin{definition}[Dynamic state]
\label{def:taxonomy-dynamic-state}
A dynamic state is a runtime state participating in operational state
evolution.
\end{definition}

Dynamic states may evolve according to
\[
x'(s) = f(x,u).
\]

% ============================================================
\section{Optimization and Relative States}
% ============================================================

\begin{definition}[Optimization state]
\label{def:taxonomy-optimization-state}
An optimization state is the state representation used as the object of
an optimization problem.
\end{definition}

Within AIM, optimization states are preferably realized runtime-state
trajectories rather than isolated geometric parameters, for example
\[
x_{\mathrm{opt}}(s) = x_{\mathrm{real}}(s).
\]

\begin{definition}[Relative state]
\label{def:taxonomy-relative-state}
A relative state is a state evaluated relative to another runtime state
or runtime reference frame.
\end{definition}

Many operational quantities are relative states, including vehicle
states relative to track states and balance states relative to
admissibility envelopes.

% ============================================================
\section{State Transformations and Hierarchy}
% ============================================================

Different state classes are connected through runtime operators. The
principal operator structure is introduced in
\cref{chap:runtime-operators}. Examples include
\[
\mathcal R,
\qquad
\mathcal D,
\qquad
\mathcal O,
\qquad
\mathcal V.
\]

A canonical hierarchy can be read as
\[
\kappa(s)
\rightarrow
\mathrm{pose2}
\rightarrow
\mathrm{pose3}
\rightarrow
x(s)
\rightarrow
x_{\mathrm{vehicle}}
\rightarrow
x_{\mathrm{CG}}
\rightarrow
x_{\mathrm{opt}}.
\]

Realization operators act throughout this hierarchy:
\[
\mathcal R :
x
\mapsto
x_{\mathrm{real}}.
\]

This hierarchy is a conceptual ordering of interpretation layers, not a
rigid implementation sequence.

% figures/foundations/state_hierarchy.tex

\begin{figure}[ht]
\centering
\begin{tikzpicture}[
  node distance=0.85cm,
  box/.style={
    draw,
    rounded corners,
    align=center,
    minimum width=4.8cm,
    minimum height=0.7cm
  },
  arrow/.style={->, thick}
]
\node[box] (kappa) {Curvature Law\\$\kappa(s)$};
\node[box, below=of kappa] (pose2) {pose2\\intrinsic planar state};
\node[box, below=of pose2] (pose3) {pose3\\railway runtime state};
\node[box, below=of pose3] (runtime) {Operational Runtime State\\$x(s)$};
\node[box, below=of runtime] (vehicle) {Vehicle State\\$x_{\mathrm{vehicle}}$};
\node[box, below=of vehicle] (cg) {CG State\\$x_{\mathrm{CG}}$};
\node[box, below=of cg] (opt) {Optimization State\\$x_{\mathrm{opt}}$};

\draw[arrow] (kappa) -- (pose2);
\draw[arrow] (pose2) -- (pose3);
\draw[arrow] (pose3) -- (runtime);
\draw[arrow] (runtime) -- (vehicle);
\draw[arrow] (vehicle) -- (cg);
\draw[arrow] (cg) -- (opt);
\end{tikzpicture}

\caption{AIM state hierarchy from intrinsic curvature to optimization state.}
\label{fig:state-hierarchy}
\end{figure}


% ============================================================
\section{Canonical Interpretation}
% ============================================================

\begin{proposition}
Within AIM, state concepts belong to distinct interpretation layers.
\end{proposition}

\begin{proposition}
pose3 describes the railway runtime state rather than the vehicle
itself.
\end{proposition}

\begin{proposition}
Vehicle states are relative runtime interpretation states generated from
interaction with the railway runtime state.
\end{proposition}

\begin{proposition}
Operational and dynamic states are generated through runtime evaluation
rather than through static geometric storage.
\end{proposition}

\begin{proposition}
Optimization states are preferably realized runtime-state trajectories
rather than isolated geometric parameters.
\end{proposition}

% ============================================================
\section{Connection to AIM}
% ============================================================

\begin{summary}
State taxonomy provides the structural interpretation framework for AIM
runtime architecture.

Within this framework, pose2 represents intrinsic reconstruction state,
pose3 represents railway runtime state, realized states represent
physical railway realization, operational states describe runtime
behaviour, vehicle states describe relative operational interpretation,
and optimization states describe the objects of trajectory optimization.

This taxonomy prevents conceptual ambiguity between geometry, runtime,
realization, dynamics, vehicle interpretation, and optimization.
\end{summary}
